\documentclass[11pt]{article}

\usepackage[T1]{fontenc}
\usepackage[utf8]{inputenc}
\usepackage{lmodern}
\usepackage[a4paper,margin=28mm]{geometry}
\usepackage{amsmath,amssymb,amsthm,mathtools,mathrsfs}
\usepackage{microtype}
\usepackage{enumitem}
\usepackage{booktabs}
\usepackage{xcolor}
\usepackage[colorlinks=true,linkcolor=blue!45!black,citecolor=blue!45!black,
  urlcolor=blue!55!black]{hyperref}
\setlength{\emergencystretch}{2em}

\newtheorem{theorem}{Theorem}[section]
\newtheorem{proposition}[theorem]{Proposition}
\newtheorem{lemma}[theorem]{Lemma}
\newtheorem{corollary}[theorem]{Corollary}
\newtheorem*{maintheorem}{Orientation-before-completion theorem}
\theoremstyle{definition}
\newtheorem{definition}[theorem]{Definition}
\theoremstyle{remark}
\newtheorem{remark}[theorem]{Remark}

\newcommand{\C}{\mathbb C}
\newcommand{\R}{\mathbb R}
\newcommand{\N}{\mathbb N}
\newcommand{\Z}{\mathbb Z}
\newcommand{\Res}{\operatorname*{Res}}
\newcommand{\ord}{\operatorname{ord}}
\newcommand{\Rea}{\operatorname{Re}}
\newcommand{\singsupp}{\operatorname{sing\,supp}}
\newcommand{\dd}{\,\mathrm d}

\title{\textbf{The admission--separation law}\\
\large Source orientation and the completion geometry of the Riemann divisor}
\author{Michael Haufschild}
\date{}

\begin{document}
\maketitle

\begin{abstract}
A singular, unbounded, or asymptotic completion is a partial quotient: it first selects an
admissible source domain and then identifies source states through an observation map.  We prove
the admission--separation law.  A target is globally recoverable precisely when every relevant
source state is admitted and the observation kernel on that domain is contained in the target
kernel.  Postprocessing changes neither gate.  For an involution, an admitted reflection pair
with nonzero odd charge must either remain separated or lose that charge.  Before completion the
two sectors are regular and complementary:
\[
 v=\frac{v+Cv}{2}+\frac{v-Cv}{2};
\]
singularity appears only when the discarded anti-invariant summand is reconstructed from the
invariant one.  Dense fibres and all-orders-flat filtrations give the corresponding continuous
and asymptotic impossibility theorems.

For the completed Riemann divisor, reflection is
$\rho\mapsto1-\overline\rho$ and the odd coordinate is
$c_\rho=\Rea\rho-\tfrac12$.  We identify both gates by source-defined objects.  The entire
character $M_\rho(z)=\exp(z(\rho-\tfrac12))$ separates every noncritical reflected pair and
recovers $c_\rho$ from one real-ray modulus.  Independently, the source-normal Hedenmalm prime row
has the exact Wick decomposition
\[
 E_N(z)=V_N(\Rea z)+R_N(\Ima z),
\]
and its fixed sub-$\log^2$ domain is classified exactly:
\[
 E_N(z)=o((\log N)^2)\quad\Longleftrightarrow\quad\Ima z=0.
\]
Thus the completion domain itself is precisely the centered critical axis.  A literal theta atom
provides a nonvanishing Gamma--Stokes source factor; Euler multiplication kills it exactly on the
Xi divisor and preserves multiplicity.  Translated Xi currents give an infinite,
multiplicity-exact singular carrier.

Explicit functional-equation-preserving quartets retain fixed nonzero charge while escaping every
compact jet, bounded Ward window, reciprocal-square energy, and real-aperture differential
matroid.  Reflected prime-valuation sectors cannot be sewn by any finite-order Rees operator.
Consequently no regular post-completion construction can supply a missing admission estimate or
recreate erased orientation.  These results do not prove the Riemann hypothesis: the remaining
statement is the source-domain theorem that every Xi root block is admitted to the fixed
prime-normal completion.  The law proves why this source theorem, rather than another downstream
positivity or spectral model, is the indispensable missing input.
\end{abstract}

\medskip
\noindent\textbf{2020 Mathematics Subject Classification.}
Primary 11M26, 11M41; Secondary 46B99, 60E15, 12H05.

\noindent\textbf{Keywords.}
Riemann hypothesis; Riemann xi function; partial completion; reflection orientation; Hedenmalm;
Gamma factor; singular support; Rees filtration.

\section{Introduction}

The functional equation of the completed zeta function preserves the divisor but reverses its
horizontal orientation.  This elementary distinction is repeatedly lost in spectral approaches
to the Riemann hypothesis.  A positive form may retain a divisor while forgetting side; a
squared normal coordinate retains location while forgetting direction; a compact observer may
see every fixed jet while a divisor packet escapes to infinite height.  Hilbert--P\'olya,
trace-formula, and theta-spectral programs therefore face the same prior question: at what stage
must the sign of $\Rea\rho-\tfrac12$ enter?  See
\cite{Connes1999,Meyer2005,ConnesConsani2021,Hedenmalm2026} for representative spectral and
operator frameworks.

Answer has two logically prior parts.  A completion may fail to admit a source state at all, as
for an unbounded graph norm or a divergent all-prime energy.  Among admitted states, every
observation $f:D\to Y$ erases its kernel pair
\[
 K_f=\{(x,y)\in D^2:f(x)=f(y)\}.
\]
A downstream operation cannot enlarge $D$ and can only enlarge $K_f$.  Exact global recovery
therefore requires universal source admission and target separation on the admitted domain.
A collection of source channels repairs the second gate exactly when their joint kernel separates
every target-relevant collision.  This is the admission--separation law proved in
Section~\ref{sec:universal}.

Reflection turns the second gate into an exact trilemma.  Write $H=1+\varepsilon$ for the
symmetric channel and $A=1-\varepsilon$ for orientation.  The identity $AH=0$ says that the
symmetric channel kills the anti-invariant summand.  It does not make source orientation
intrinsically singular: when $2$ is invertible, the invariant and anti-invariant summands are
complementary and reconstruct the source.  For nonzero odd charge, a completion must retain that
second summand, identify the reflected pair and lose the charge, or exclude a member of the pair
from its regular domain.  Topological completion converts failed separation into dense-fibre
collapse; asymptotic completion converts failed admission or sewing into a beyond-all-orders
cost.

For a Riemann zero $\rho$, the odd target is
\[
 c_\rho=\Rea\rho-\frac12,
 \qquad c_{1-\overline\rho}=-c_\rho.
\]
The centered character $M_\rho$ is the minimal scalar source channel retaining this charge.
Its modulus on any nonzero real ray recovers $c_\rho$ exactly.  This character is not a decorative
reparametrization.  It is simultaneously the cofinal shell rate, the principal Hurwitz phase,
and the transverse factor multiplying the local divisor residue.  Section~\ref{sec:riemann}
proves these identities.

Three exact source theorems then locate the two gates.  First, the normal multiplier of the
Hedenmalm theta-interval row decomposes termwise into vertical trigonometric and radial hyperbolic
prime energies.  The vertical term is at most harmonic while every radial displacement pays the
quadratic prime susceptibility.  The resulting sub-$\log^2$ domain is exactly the centered real
axis.  Second, the normalized theta seed has Mellin
transform equal to the archimedean factor.  Its causal and anti-causal Green modes differ by a
nonvanishing homogeneous Gamma character; multiplying by the Euler factor annihilates this jump
exactly on the Xi divisor and does not alter multiplicity.  Third, the translated logarithmic
Xi current has a removable residue germ whose value is the finite difference of translated
multiplicities.  Every positive translation has infinitely many such jumps.  Thus the
orientation multiplier is necessarily singular on an infinite carrier, not merely at an
exceptional point.

Finally, two sharp no-go results show that regularization cannot recover what completion erased.
Explicit off-line quartets can escape every compact jet and bounded Ward window without changing
their modular charge.  Prime-valuation reflection produces an all-orders-flat sector paired with
a finite-grade sector; every exact sewing has essential norm growth.  Together these results
explain the dense radicals, Wold--Poisson transport, Sonin--Calkin seams, Ward charges, and
finite-reporter failures that otherwise look route-specific.  The law does not create the
one-sided arithmetic inequality needed for zero exclusion.  It proves where such orientation
must live and which broad classes of post-completion repair cannot contain it.

The admission--separation factorization and algebraic splitting are elementary; no novelty is
claimed for them in isolation.  Contribution is their exact Riemann realization: one fixed
source-normal completion whose regular domain is precisely the centered critical axis, one
canonical character which separates reflected orientation, and sharp proofs that downstream
compact or finite-grade repairs cannot replace either gate.  We do not assert that every Xi root
block enters the fixed domain.  That source-admission assertion, not universal unitarity of all
characters, is the remaining RH-bearing theorem.

\subsection*{Main theorem}

\begin{maintheorem}\label{thm:main}
The following statements hold unconditionally.
\begin{enumerate}[label=\textup{(\roman*)}]
\item Let a partial completion be a pair $(D,\pi)$ with $D\subseteq X$ and
$\pi:D\to Q$.  It recovers $\chi:X\to C$ globally exactly when
\[
 D=X\qquad\hbox{and}\qquad K_\pi\subseteq K_{\chi|D}.
\]
Postprocessing preserves $D$ and enlarges $K_\pi$.  More generally, for any heterogeneous
family $\tau_j:D\to T_j$, the target factors uniquely on the attained joint range exactly when
\[
 K_\pi\cap\bigcap_jK_{\tau_j}\subseteq K_{\chi|D}.
\]

If $D$ is stable under an involution $\sigma$, $\pi\sigma=\pi$, and
$\chi\sigma=-\chi$ with $\chi$ real-valued, then $\chi$ cannot factor through $\pi$ unless
$\chi=0$ on $D$.  Every sufficient transverse channel separates $x$ from $\sigma x$ wherever
$\chi(x)\ne0$.  In a commutative ring, $\varepsilon^2=1$ implies
$(1-\varepsilon)(1+\varepsilon)=0$; if $2$ is regular, then
\[
 \varepsilon\ne1\quad\Longleftrightarrow\quad1+\varepsilon
 \text{ is not regular}.
\]
If $C$ is a linear involution and $2$ is invertible, its invariant and anti-invariant modules are
complementary and every vector is recovered by
\[
 v=\frac{v+Cv}{2}+\frac{v-Cv}{2}.
\]

If the fibres of $\pi$ are dense, a $T_1$-valued downstream reporter whose pullback is
continuous at one point is constant.  If one collapsed source germ is flat to every order of a
scale $a\to0$ while the target has a nonzero finite grade, every exact linear sewing norm
dominates $a^{-K}$ for every $K$.

\item Define the source-normal prime energy
\[
 E_N(z)=\sum_{p\le N}\frac{4}{p}
   \left|\sin\!\left(\frac{z\log p}{2}\right)\right|^2.
\]
Writing $z=x+iy$, it decomposes exactly as
\[
 E_N(z)=4\sum_{p\le N}\frac{\sin^2(x\log p/2)}p
       +4\sum_{p\le N}\frac{\sinh^2(y\log p/2)}p.
\]
Its sub-$\log^2$ domain is precisely the centered real axis:
\[
 \frac{E_N(z)}{(\log N)^2}\longrightarrow0
 \quad\Longleftrightarrow\quad \Ima z=0.
\]

Let $\rho$ lie in the open critical strip and put
$M_\rho(z)=e^{z(\rho-1/2)}$.  Critical-line reflection replaces $c_\rho$ by $-c_\rho$ and
replaces real-ray character norms by their reciprocals.  Any nonzero real ray recovers
$c_\rho$ exactly.  At a zero of $\xi$ of multiplicity $m$, the centered outer current
\[
 \mathcal J_\xi(s)=-\left(s-\frac12\right)\frac{\xi'(s)}{\xi(s)}
\]
has residue $-m(\rho-\tfrac12)$.  The cofinal shell modulus and principal Hurwitz phase recover
the same horizontal charge.

The Gamma--Stokes jump of the literal theta seed is nonzero throughout the strip, is killed by
Euler multiplication exactly at Xi zeros, and preserves analytic multiplicity.  Every positive
translation of the Xi logarithmic current has infinitely many multiplicity jumps, each lying in
the singular support of the canonical orientation multiplier.

\item For every $0<c<1/2$ there are functional-equation-preserving hypertranscendental quartet
perturbations with fixed charge $c$ which converge with every holomorphic jet on compact sets,
whose Ward density converges uniformly to zero on every bounded window, whose reciprocal-square
energy tends to zero, and whose complete real-aperture differential matroid is unchanged.  The
modular character still recovers $c$, and total Ward mass remains $8\pi$ per quartet.

In the reflected prime-valuation realization, the positive-charge squared affinity is flat beyond
every power of the Wold boundary volume while the negative-charge squared affinity has finite
Wold grade (and its unsquared affinity has half-normal grade).  No finite-Rees-moderate operator
sews the pair.  Every exact sewing operator $R_N$
satisfies
\[
 a_N^{-K}=o(\|R_N\|)\qquad(K\in\N),
 \quad a_N=\prod_{p<N}(1-p^{-1}).
\]
\end{enumerate}
\end{maintheorem}

\section{The universal law}\label{sec:universal}

\subsection{Completion is a partial quotient}

A partial completion of a source set $X$ is a pair $(D,\pi)$ consisting of an admission domain
$D\subseteq X$ and an observation $\pi:D\to Q$.  A target $\chi:X\to C$ is globally recovered
when $D=X$ and there is a function on the attained range $\pi(D)$ whose pullback is $\chi$.

\begin{theorem}[Admission--separation law]\label{thm:admission}
The partial completion $(D,\pi)$ globally recovers $\chi$ if and only if
\[
 D=X\quad\hbox{and}\quad K_\pi\subseteq K_{\chi|D}.
\tag{2.1}\label{eq:admission-separation}
\]
For every postprocessing map $g:Q\to Y$, the partial completion $(D,g\circ\pi)$ has the same
admission domain and $K_\pi\subseteq K_{g\circ\pi}$.
\end{theorem}

\begin{proof}
Global recovery includes $D=X$ by definition.  On $D$, factorization through the attained range
is equivalent to constancy of $\chi$ on every $\pi$-fibre, which is the kernel inclusion.
Postprocessing neither changes $D$ nor separates two equal $\pi$-values.
\end{proof}

The theorem separates two failures which a total-map model conflates.  A source point outside
$D$ has no completed value; two points inside $D$ may nevertheless have the same value.  No
downstream operation repairs either failure.  Domain estimates and transverse reporters are
therefore complementary obligations, not alternative proof strategies.

\subsection{Kernel pairs and exact recovery}

For a map $f:X\to Y$, write $K_f(x,y)$ for $f(x)=f(y)$.  For a heterogeneous family
$\tau_j:X\to T_j$, define its record
\[
 \tau(x)=(\tau_j(x))_j\in\prod_jT_j.
\]

\begin{theorem}[Kernel-pair recovery law]\label{thm:kernel}
Let $f:X\to Y$ and $\chi:X\to C$.  There is a function
$d:f(X)\to C$ such that $\chi=d\circ f$ if and only if
\[
 K_f\subseteq K_\chi.
\tag{2.2}\label{eq:kernel-refinement}
\]
Consequently, $\chi$ is recoverable from the record $(\pi,\tau)$ if and only if
\[
 K_\pi\cap\bigcap_jK_{\tau_j}\subseteq K_\chi.
\tag{2.3}\label{eq:joint-kernel}
\]
If $\chi=d\circ\pi$ on all of $Q$, then $K_\pi\subseteq K_\chi$.
\end{theorem}

\begin{proof}
If $\chi=d\circ f$, equality of $f$-values implies equality of $\chi$-values.  Conversely,
define $d(f(x))=\chi(x)$.  Condition \eqref{eq:kernel-refinement} makes this independent of the
chosen preimage.  Apply the result to the joint record; its kernel is exactly the intersection
in \eqref{eq:joint-kernel}.  Last assertion is the forward implication.
\end{proof}

The theorem makes no topological, algebraic, metric, probabilistic, boundedness, or finiteness
assumption.  It also separates necessity from construction: a new channel repairs completion
only by splitting every completion collision on which the target differs.

\begin{corollary}[Reflection orientation must be transverse]\label{cor:odd}
Let $\sigma:X\to X$ be an involution, let $\pi\sigma=\pi$, and let
$\chi:X\to\R$ satisfy $\chi(\sigma x)=-\chi(x)$.  If $\chi$ factors through $\pi$, then
$\chi=0$.  If $(\pi,\tau)$ recovers $\chi$ and $\chi(x)\ne0$, then
$\tau(\sigma x)\ne\tau(x)$.  For scalar $\tau$, its odd part
\[
 \tau^-(x)=\frac{\tau(x)-\tau(\sigma x)}2
\]
is nonzero at every such $x$.
\end{corollary}

\begin{proof}
Invariance gives $\pi(\sigma x)=\pi(x)$.  Factorization would give
$\chi(\sigma x)=\chi(x)$, while oddness gives $\chi(\sigma x)=-\chi(x)$; hence
$\chi(x)=0$.  For joint recovery, equality of the $\tau$-values would place the reflected pair
in both kernels in \eqref{eq:joint-kernel}, again forcing $\chi(x)=0$.  The scalar assertion is
equivalent to separation.
\end{proof}

\subsection{Symmetric completion kills a complementary sector}

\begin{theorem}[Orientation annihilator law]\label{thm:singular-support}
Let $A$ be a commutative ring and let $\varepsilon\in A$ satisfy $\varepsilon^2=1$.  Then
\[
 (1-\varepsilon)(1+\varepsilon)=0.
\tag{2.4}\label{eq:orientation-annihilator}
\]
If $1+\varepsilon$ is regular, then $\varepsilon=1$.  If $2$ is regular, then
\[
 \varepsilon\ne1\quad\Longleftrightarrow\quad
 1+\varepsilon\text{ is not regular}.
\tag{2.5}\label{eq:regular-equivalence}
\]
For an $A$-module $V$, every anti-invariant vector $v$ satisfying
$\varepsilon v=-v$ obeys $(1+\varepsilon)v=0$.
\end{theorem}

\begin{proof}
Identity \eqref{eq:orientation-annihilator} is $1-\varepsilon^2=0$.  Regularity cancels
$1+\varepsilon$ and gives $1-\varepsilon=0$.  Conversely, if $\varepsilon=1$, then
$1+\varepsilon=2$, regular by hypothesis.  Module statement is
$v+\varepsilon v=v-v=0$.
\end{proof}

The annihilator concerns recovery from the symmetric channel; it does not make source orientation
singular.  There is a canonical positive counterpart.

\begin{proposition}[Faithful orientation descent]\label{prop:orientation-split}
Let $V$ be a module over a commutative ring in which $2$ is invertible, and let $C:V\to V$ be a
linear involution.  The invariant and anti-invariant modules
\[
 V^+=\ker(C-1),\qquad V^-=\ker(C+1)
\]
are complementary.  Explicitly,
\[
 v=v^++v^-,\qquad
 v^+=\frac{v+Cv}{2},\qquad v^-=\frac{v-Cv}{2}.
\tag{2.6}\label{eq:orientation-split}
\]
In particular, the joint invariant--anti-invariant record is faithful.
\end{proposition}

\begin{proof}
The displayed terms are respectively fixed and negated by $C$, and their sum is $v$.  A vector
in both submodules satisfies $v=-v$, hence $2v=0$ and therefore $v=0$.
\end{proof}

Thus orientation is a regular complementary source coordinate.  It becomes torsion or singular
only when one asks the symmetric completion alone to reconstruct the summand it annihilated.

\subsection{Topological and filtered forms}

\begin{proposition}[Dense-fibre collapse]\label{prop:dense}
Let $\pi:X\to Q$ be a surjection from a topological space, suppose every fibre is dense, and let
$Y$ be $T_1$.  If $F:Q\to Y$ has $F\circ\pi$ continuous at one point, then $F$ is constant.
\end{proposition}

\begin{proof}
Suppose $F(q_1)\ne F(q_2)$ and choose $x$ over $q_1$.  A $T_1$ neighbourhood of $F(q_1)$
excluding $F(q_2)$ pulls back to a neighbourhood of $x$.  Density of the fibre over $q_2$
forces that neighbourhood to contain a point mapped to $F(q_2)$, contradiction.
\end{proof}

Let $\ell$ be a filter, let $a$ be nonzero eventually, and write
$f=o_\ell(a^m)$ in the usual asymptotic sense.

\begin{theorem}[Beyond-all-orders sewing]\label{thm:rees}
Let $f$ be eventually nonzero and satisfy $f=o_\ell(a^m)$ for every $m\in\N$.  Let
$g=\Theta_\ell(a^d)$ for some finite $d$.  Then for every $K$,
\[
 a^{-K}=o_\ell(g/f).
\tag{2.7}\label{eq:scalar-sewing}
\]
More generally, if $f(x)$ and $g(x)$ take values in real normed spaces and bounded linear maps
$R_x$ satisfy $R_xf(x)=g(x)$, then
\[
 a^{-K}=o_\ell(\|R_x\|)\qquad(K\in\N).
\tag{2.8}\label{eq:operator-sewing}
\]
\end{theorem}

\begin{proof}
Use flatness at order $K+d$: $f=o(a^{K+d})$.  Taking reciprocals reverses the little-$o$
relation, and multiplication by $g=\Theta(a^d)$ gives \eqref{eq:scalar-sewing}.  For the operator
form,
$\|g\|\le\|R_x\|\,\|f\|$, so $\|g\|/\|f\|=O(\|R_x\|)$; compose with the scalar estimate.
\end{proof}

Theorems~\ref{thm:admission}, \ref{thm:kernel}, \ref{thm:singular-support}, and \ref{thm:rees}
express one irreversible order of information.  Admission specifies where completion exists;
kernel inclusion gives algebraic recoverability; fibre closure gives continuous regularity; and
asymptotic grade quantifies the unavoidable cost of sewing incompatible domains.

\section{The Riemann orientation charge}\label{sec:riemann}

Write
\[
 \xi(s)=\frac12s(s-1)\pi^{-s/2}\Gamma(s/2)\zeta(s),
 \qquad \xi(s)=\xi(1-s),
\tag{3.1}\label{eq:xi}
\]
and let $\rho$ lie in the open critical strip.  Schwarz reflection and the functional equation
pair $\rho$ with
\[
 \sigma\rho=1-\overline\rho.
\]

\subsection{The exact source-normal admission domain}

Remove the common root coordinate from the normal multiplier of the Hedenmalm theta-interval
row.  Its prime component is
\[
 b_p(z)=2p^{-1/2}\sin\!\left(\frac{z\log p}{2}\right),
\]
so the squared energy of the first $N$ prime components is
\[
 E_N(z)=\sum_{p\le N}|b_p(z)|^2
 =\sum_{p\le N}\frac4p
   \left|\sin\!\left(\frac{z\log p}{2}\right)\right|^2.
\tag{3.2}\label{eq:prime-normal-energy}
\]

\begin{theorem}[Exact prime-normal domain]\label{thm:prime-normal-domain}
For $z=x+iy$,
\[
 E_N(z)=
 4\sum_{p\le N}\frac{\sin^2(x\log p/2)}p+
 4\sum_{p\le N}\frac{\sinh^2(y\log p/2)}p.
\tag{3.3}\label{eq:wick-split}
\]
Consequently the fixed sub-$\log^2$ completion has the exact domain
\[
 \frac{E_N(z)}{(\log N)^2}\longrightarrow0
 \quad\Longleftrightarrow\quad y=0.
\tag{3.4}\label{eq:prime-normal-domain}
\]
\end{theorem}

\begin{proof}
Identity \eqref{eq:wick-split} is the termwise formula
$|\sin(u+iv)|^2=\sin^2u+\sinh^2v$.  If $y=0$, then
\[
 0\le E_N(x)\le4\sum_{p\le N}\frac1p
 \le4\sum_{n\le N}\frac1n=O(\log N),
\]
which proves the forward domain estimate.  If $y\ne0$, the inequality
$\sinh^2v\ge v^2$ gives
\[
 E_N(x+iy)\ge
 y^2\sum_{p\le N}\frac{(\log p)^2}{p}.
\]
The prime number theorem and partial summation give
$\sum_{p\le N}(\log p)^2/p\sim\tfrac12(\log N)^2$; hence the normalized
energy cannot tend to zero.  This proves the converse.
\end{proof}

The centered spectral coordinate is $z=-i(\rho-\tfrac12)$, for which
$\Ima z=\tfrac12-\Rea\rho=-c_\rho$.  The theorem therefore classifies the regular domain itself;
it does not assert that a Xi root block belongs to it.  Universal admission of those root blocks
is exactly the still-missing RH-bearing source estimate.

\subsection{The centered modular character}

\begin{proposition}[Exact character recovery]\label{prop:character}
Define
\[
 M_\rho(z)=\exp\!\left(z(\rho-\tfrac12)\right).
\tag{3.5}\label{eq:character}
\]
Then $M_\rho(z+w)=M_\rho(z)M_\rho(w)$ and, for $t\in\R$,
\[
 |M_\rho(t)|=e^{t c_\rho},
 \qquad
 c_\rho=\frac{\log|M_\rho(t)|}{t}\quad(t\ne0).
\tag{3.6}\label{eq:real-ray-recovery}
\]
Moreover,
\[
 c_{\sigma\rho}=-c_\rho,
 \qquad |M_{\sigma\rho}(t)|\,|M_\rho(t)|=1.
\tag{3.7}\label{eq:character-reflection}
\]
For $0<\Rea\rho<1$, the principal phase on the Hurwitz period satisfies
\[
 \frac{\left|\arg M_\rho(-2\pi i)\right|}{2\pi}=|c_\rho|.
\tag{3.8}\label{eq:hurwitz-phase}
\]
\end{proposition}

\begin{proof}
All identities follow from the exponential law and
$\Rea(t(\rho-1/2))=tc_\rho$.  Reflection negates $c_\rho$.  Finally
$M_\rho(-2\pi i)=-e^{-2\pi i\rho}$; the open-strip bound
$|c_\rho|<1/2$ places its centered argument in the principal branch and gives
\eqref{eq:hurwitz-phase}.
\end{proof}

Thus RH is equivalent to unitarity of every divisor character on the real ray.  Proposition
\ref{prop:character} does not prove that unitarity; it identifies the exact transverse datum a
source argument must control.

For any scale $q>1$, the cofinal shell character
\[
 T_q(\rho)=q^{1/2-\rho}
\]
has
\[
 \frac{\log|T_q(\rho)|}{\log q}=-c_\rho.
\tag{3.9}\label{eq:shell-rate}
\]
The modular, Hurwitz, and shell charts therefore carry the same charge.

\subsection{The centered outer current}

\begin{proposition}[Multiplicity-weighted outer charge]\label{prop:outer-current}
The meromorphic current
\[
 \mathcal J_\xi(s)=-\left(s-\frac12\right)\frac{\xi'(s)}{\xi(s)}
\tag{3.10}\label{eq:outer-current}
\]
is invariant under $s\mapsto1-s$.  If $\rho$ is a zero of multiplicity $m$, then
\[
 \Res_{s=\rho}\mathcal J_\xi(s)=-m(\rho-\tfrac12).
\tag{3.11}\label{eq:outer-residue}
\]
Its real part is $-mc_\rho$.
\end{proposition}

\begin{proof}
Differentiating $\xi(1-s)=\xi(s)$ shows that $\xi'/\xi$ is reflection-odd; the affine factor is
also odd, so their product is invariant.  Locally
$\xi(s)=(s-\rho)^m u(s)$ with $u(\rho)\ne0$, hence
$\xi'/\xi=m/(s-\rho)+u'/u$.  Equation \eqref{eq:outer-residue} follows.
\end{proof}

This residue is the outer connection of the centered character: multiplicity is longitudinal,
while $\rho-1/2$ is the transverse source charge.  Squaring that coordinate preserves the
zero set but destroys its reflection orientation.

\section{The singular source carrier}\label{sec:source}

\subsection{A literal Gamma--Stokes factor}

Let
\[
 k(t)=2\pi(2\pi t^4-3t^2)e^{-\pi t^2},\qquad t>0,
\tag{4.1}\label{eq:theta-seed}
\]
and put
\[
 A(s)=\frac{s(s-1)}2\pi^{-s/2}\Gamma(s/2).
\tag{4.2}\label{eq:arch-factor}
\]
Direct Mellin integration gives
\[
 \int_0^\infty t^{s-1}k(t)\,\dd t=A(s),\qquad \Rea s>0.
\tag{4.3}\label{eq:mellin-seed}
\]
Define causal and anti-causal modes
\[
 C_s(t)=t^{-s}\int_0^t u^{s-1}k(u)\,\dd u,
 \qquad
 A_s(t)=-t^{-s}\int_t^\infty u^{s-1}k(u)\,\dd u.
\tag{4.4}\label{eq:stokes-modes}
\]

\begin{theorem}[Gamma--Stokes divisor factorization]\label{thm:gamma-stokes}
For $\Rea s>0$ and $t>0$,
\[
 C_s(t)-A_s(t)=t^{-s}A(s).
\tag{4.5}\label{eq:stokes-jump}
\]
In the open critical strip this jump is nonzero and
\[
 \zeta(s)\bigl(C_s(t)-A_s(t)\bigr)=0
 \quad\Longleftrightarrow\quad \xi(s)=0.
\tag{4.6}\label{eq:euler-kills}
\]
Multiplication by the nonvanishing character $e^{-s\log t}$ does not change the analytic
multiplicity of a zero of $\xi$.
\end{theorem}

\begin{proof}
Splitting the Mellin integral at $t$ gives \eqref{eq:stokes-jump}.  In the strip, $t^{-s}$ and
$A(s)$ are nonzero.  Since $\xi(s)=A(s)\zeta(s)$ there, \eqref{eq:euler-kills} follows.
The last assertion follows because a nonvanishing analytic factor has order zero.
\end{proof}

The theorem supplies the source-pointed object missing from a regular completed
observer.  The homogeneous character exists before the Euler factor selects the divisor; Euler
annihilation then records the divisor without changing its multiplicity.

\subsection{Multiplicity transgression and infinite singular support}

Let $m_\xi(z)$ denote analytic zero multiplicity, with value zero away from the divisor, and let
\[
 G_h(s)=\frac{\xi'(s)}{\xi(s)}-\frac{\xi'(s-h)}{\xi(s-h)}.
\tag{4.7}\label{eq:shift-current}
\]
For $a>0$, define the Hedenmalm frequency coordinate and horizontal orientation multiplier
\[
 b_a(z)=-i\bigl(z-(\tfrac12+a)\bigr),\qquad
 \mathcal Q_a(b)=
 \frac{\xi'}{\xi}(\tfrac12+a+ib)+
 \frac{\xi'}{\xi}(\tfrac12+a-ib).
\tag{4.8}\label{eq:hedenmalm-multiplier}
\]
The functional equation gives
$G_{2a}(z)=\mathcal Q_a(b_a(z))$.

\begin{theorem}[Infinite divisor transgression]\label{thm:transgression}
For every $h>0$, the jump set
\[
 \mathcal D_h=\{\rho\in\C:m_\xi(\rho-h)\ne m_\xi(\rho)\}
\tag{4.9}\label{eq:jump-set}
\]
is infinite.  At each $\rho\in\mathcal D_h$, $G_h$ is nonanalytic, and the removable residue
germ
\[
 x\longmapsto xG_h(\rho+x)
\]
has center value
\[
 m_\xi(\rho)-m_\xi(\rho-h)
\tag{4.10}\label{eq:multiplicity-jump}
\].
Under the affine coordinate $b_a$ with $h=2a$, these points give an infinite singular subset of
$\mathcal Q_a$ at every $a>0$.
\end{theorem}

\begin{proof}
Local logarithmic-derivative expansions give \eqref{eq:multiplicity-jump}; a nonzero residue
prevents analyticity.  It remains to prove infinitude.  Translation by $h$ partitions the Xi
divisor into horizontal $h\Z$-orbits.  Each orbit contains only finitely many zeros because all
nontrivial zeros lie in $0<\Rea s<1$.  If only finitely many orbits had a left boundary, the
divisor would be finite, contradicting the classical infinitude of Xi zeros.  Hence infinitely
many orbit-boundary points satisfy $m_\xi(\rho)>0=m_\xi(\rho-h)$ and lie in $\mathcal D_h$.
Affine frequency change is injective and preserves infinitude and nonanalyticity.
\end{proof}

Theorem~\ref{thm:transgression} turns the abstract statement ``orientation lives on singular
support'' into a multiplicity-exact Riemann theorem.  The carrier is infinite at every positive
aperture; no simplicity hypothesis is used.

\section{Sharp failure of regular recovery}\label{sec:nogo}

\subsection{Compact and differential blindness}

We first record why the completed function itself has maximal differential complexity.
Put $L=\xi'/\xi$.  Gamma recurrence and the Euler series give, on $\Rea s>1$,
\[
 L(s+2)-L(s)-\left(\frac1{s+2}+\frac1{s+1}-\frac1{s-1}\right)
 =\sum_{n\ge1}\frac{\Lambda(n)(1-n^{-2})}{n^s}.
\tag{5.1}\label{eq:xi-shift-defect}
\]
The coefficient at every prime is nonzero.  Ostrowski's prime-support criterion therefore makes
the right side hypertranscendental over $\C(s)$ \cite{Ostrowski1920,Nagoshi2016}.  If $\xi$ were
differentially algebraic, then so would be $L$, its translate, and the left side, a contradiction.
Hence $\xi$ is hypertranscendental over $\C(s)$.

Fix $0<c<1/2$, put $z=s-1/2$, and for $B>0$ define
\[
 Q_{c,B}(s)=\left(1+\frac{(z-c)^2}{B^2}\right)
             \left(1+\frac{(z+c)^2}{B^2}\right).
\tag{5.2}\label{eq:quartet}
\]
Its zeros are $1/2\pm c\pm iB$, it satisfies the functional equation and Schwarz symmetry, and
$Q_{c,B}\to1$ with every derivative on compact sets as $B\to\infty$.  Put
$F_{N,B}=Q_{c,B}^N\xi$.

For a nonzero meromorphic function $F$ and real aperture $h$, write
\[
 D_h^F(s)=\frac{F'(s+h)}{F(s+h)}-\frac{F'(s)}{F(s)}.
\tag{5.3}\label{eq:aperture-family}
\]
For a family in a differential extension of $\C(s)$, its represented differential matroid records
which finite subfamilies are differentially algebraically independent over $\C(s)$.

\begin{theorem}[Escaping orientation charge]\label{thm:escape}
For every $N\ge1$, $F_{N,B}$ is hypertranscendental, has an inserted
off-critical quartet of multiplicity $N$, and converges to $\xi$ with all holomorphic jets on
compact sets as $B\to\infty$.  The complete real-aperture differential matroid
$\operatorname{Mat}(D_h^{F_{N,B}})_{h\in\R}$ equals
$\operatorname{Mat}(D_h^\xi)_{h\in\R}$.  Its reciprocal-square orientation energy
\[
 E_{c,B}=\frac{\sin^2(2\pi c)}{c^2+B^2}
\tag{5.4}\label{eq:trace-energy}
\]
tends to zero.  The associated nonnegative Ward density is
\[
 W_{c,B}^{(N)}(t)=2N\left(
 \frac{2(c+1/2)}{(c+1/2)^2+(t+B)^2}
 +\frac{2(c+1/2)}{(c+1/2)^2+(t-B)^2}
 \right),
\tag{5.5}\label{eq:ward-density}
\]
and satisfies
\[
 \int_\R W_{c,B}^{(N)}(t)\,\dd t=8\pi N,
 \qquad
 W_{c,B}^{(N)}\longrightarrow0
 \text{ uniformly on bounded windows}.
\tag{5.6}\label{eq:ward-escape}
\]
The centered modular character of every inserted point still recovers the fixed charge $c$.
\end{theorem}

\begin{proof}
All divisor symmetries follow from \eqref{eq:quartet}.  Its expansion is
$1+2(z^2+c^2)B^{-2}+(z^2-c^2)^2B^{-4}$, proving all-jet compact convergence; multiplication by a
nonzero polynomial preserves hypertranscendence.  If $q_{c,B}=Q'_{c,B}/Q_{c,B}$, then
\[
 D_h^{F_{N,B}}=D_h^\xi+N\bigl(q_{c,B}(s+h)-q_{c,B}(s)\bigr).
\]
The correction is rational, so translation by it does not change differential dependence over
$\C(s)$.  Formula \eqref{eq:trace-energy} gives reciprocal-square decay.  Each term in
\eqref{eq:ward-density} is a
positive translated Lorentzian.  The identity
$\int_\R 2a/(a^2+t^2)\,\dd t=2\pi$ gives the mass.  On $|t|\le R<B$, both denominators are at
least $(B-R)^2$, giving uniform decay.  Character recovery is
\eqref{eq:real-ray-recovery}.
\end{proof}

The theorem proves simultaneous blindness of compact analysis, differential algebra, bounded
Ward observation, and trace-class decay.  It also identifies the positive survivor: a
height-uniform source character.  The total Ward index counts off-line defects; it does not make
their number vanish.

\subsection{Prime-valuation type fracture}

For a prime $p$ and $|c|<1/2$, let
\[
 q_p=p^{-1},\qquad r_{p,c}=p^{-(1-2c)},\qquad u_{p,c}=p^{-(1-c)}
\]
and define the squared Hellinger factor
\[
 h_{p,c}=\frac{(1-q_p)(1-r_{p,c})}{(1-u_{p,c})^2}.
\tag{5.7}\label{eq:hellinger-factor}
\]
Put
\[
 H_N(c)=\prod_{p<N}h_{p,c},
 \qquad a_N=\prod_{p<N}(1-p^{-1}).
\tag{5.8}\label{eq:hellinger-product}
\]
Call $R_N$ finite-Rees-moderate when $\|R_N\|=O(a_N^{-K})$ for some fixed $K$.

\begin{theorem}[Reflected asymptotic-type fracture]\label{thm:fracture}
For $0<c<1/2$,
\[
 H_N(c)=o(a_N^K)\quad\text{for every }K,
 \qquad
 H_N(-c)=\Theta(a_N).
\tag{5.9}\label{eq:type-fracture}
\]
Consequently no operator family of fixed finite Rees order sends $H_N(c)$ exactly to
$H_N(-c)$.  Every exact bounded linear sewing $R_NH_N(c)=H_N(-c)$ satisfies
\[
 a_N^{-K}=o(\|R_N\|)\qquad(K\in\N).
\tag{5.10}\label{eq:essential-cost}
\]
\end{theorem}

\begin{proof}
For $c<0$, Euler-product algebra gives
\[
 \frac{H_N(c)}{a_N}
 \longrightarrow
 \frac{\zeta(1-c)^2}{\zeta(1-2c)}>0,
\]
which is the second relation in \eqref{eq:type-fracture}.  For $c>0$, the local ratio
$r_{p,c}/q_p=p^{2c}$ is unbounded, while $u_{p,c}/r_{p,c}=p^{-c}\to0$.
Fix $L$.  All sufficiently large primes therefore satisfy
$8u_{p,c}\le3r_{p,c}$ and $4Lq_p\le r_{p,c}$.  Using
$u_{p,c}^2=q_pr_{p,c}$, direct expansion gives
\[
 h_{p,c}\le1-\frac{r_{p,c}}4
 \le1-Lq_p\le(1-q_p)^L,
\]
where the last step is Bernoulli's inequality.  Taking $L=K+1$, a finite prefix contributes only a
constant, and Mertens divergence of $\sum_p1/p$ yields
$H_N(c)=O(a_N^{K+1})=o(a_N^K)$.  Apply Theorem~\ref{thm:rees} with $d=1$ to obtain
\eqref{eq:essential-cost}; finite-Rees maps preserve all-orders flatness and therefore cannot
reach the negative sector.
\end{proof}

Reflection has therefore produced more than two different finite grades.  One side is finite
and half-normal after taking the unsquared affinity; the other is flat beyond every Rees order.
No bounded graded completion can sew the reflected pair.

\section{Consequences and scope}\label{sec:scope}

Finite reporters, tail quotients, and recurrence coronas are direct kernel-pair instances whenever
their defining equivalence erases the reported coordinates.  Reflection-even Weil forms,
positive Grams, and squared Julia coordinates are involutive instances whenever they identify
$\rho$ with $1-\overline\rho$: by Corollary~\ref{cor:odd}, none can itself carry the odd charge
$c_\rho$.  These conclusions need no spectral estimate beyond identification of the relevant
kernel.

Dense Hedenmalm radicals are the topological realization.  Once positivity and translation
covariance place a cyclic theta orbit in a form radical, ordinary-$L^2$ density makes that radical
dense.  A nonzero form then has a proper nonclosed radical and cannot be repaired by a regular
Hilbert completion.  This uses the cyclicity and positivity hypotheses of the theta problem; it is
not a set-theoretic corollary with those analytic inputs removed.

Wold--Poisson transport, Sonin--Calkin seams, and Ward energies occupy the other side of the law:
they are downstream consumers of a surviving transverse character.  Their transport and
positivity estimates do not manufacture orientation.  Theorem~\ref{thm:gamma-stokes} and
Theorem~\ref{thm:transgression} identify the upstream source and singular carrier that such
consumers require.  Theorem~\ref{thm:fracture} quantifies the price of trying to replace that
carrier by finite-grade sewing.

The remaining RH wall is now exact.  Every off-line divisor character is nonunitary on the real
source ray.  Theorems~\ref{thm:gamma-stokes} and \ref{thm:transgression} show where that character
enters and where its multiplicity lives.  Theorems~\ref{thm:escape} and \ref{thm:fracture} prove
that compact, invariant, finite-grade repair cannot remove it.  Zero exclusion therefore requires
a source-derived one-sided law that forces $c_\rho=0$ before reflection-even completion.  No such
inequality is asserted here, and no equivalence is presented as its proof.

This is a closed classification, not an open-question list:
\begin{enumerate}[label=\textup{(\alph*)}]
\item regular invariant postprocessing cannot recover Riemann orientation;
\item every sufficient repair has a nonzero odd source component off the critical line;
\item the canonical analytic carrier is necessarily singular on an infinite multiplicity-jump
set;
\item finite-grade sewing of reflected valuation sectors is impossible and exact sewing is
beyond all orders;
\item local compact and differential-algebraic agreement does not control the modular charge.
\end{enumerate}

\section{Conclusion}

Completion has a universal information order.  Its kernel pair records every distinction it has
destroyed, and no downstream operation can shrink that kernel.  Under reflection, the statement
becomes an algebraic support law: orientation is annihilated by the symmetric channel and can
survive only where that channel is singular.  Dense fibres and all-orders-flat germs are the
topological and quantitative forms of the same phenomenon.

For the Riemann divisor, the lost distinction is exactly
$c_\rho=\Rea\rho-1/2$.  The centered modular character recovers it on any nonzero real source ray;
shell transport, Hurwitz phase, and the centered divisor residue carry the same charge.  The
literal Gamma--Stokes jump supplies a nonvanishing source-pointed factor before Euler
annihilation, while translated Xi currents place multiplicity on an infinite singular carrier.
Escaping quartets and reflected valuation sectors prove that neither compact agreement nor
finite-grade operator repair can replace that carrier.

The resulting principle is simple: orient before completing.  This does not prove RH, because it
does not supply the upstream arithmetic sign forcing every centered charge to vanish.  It does
identify the mathematical object such a proof must control and closes the broad class of regular
post-completion mechanisms that attempt to recover orientation after it has already been erased.

\begin{thebibliography}{99}

\bibitem{Connes1999}
A.~Connes,
\newblock Trace formula in noncommutative geometry and the zeros of the Riemann zeta function,
\newblock \emph{Selecta Math. (N.S.)} \textbf{5} (1999), 29--106.

\bibitem{ConnesConsani2021}
A.~Connes and C.~Consani,
\newblock The scaling Hamiltonian,
\newblock \emph{J. Operator Theory} \textbf{85} (2021), 257--276.

\bibitem{GauthierXarles2009}
P.~M. Gauthier and X.~Xarles,
\newblock Perturbations of $L$-functions with or without non-trivial zeros off the critical line,
\newblock arXiv:0911.5135 (2009).

\bibitem{Hedenmalm2026}
H.~Hedenmalm,
\newblock Spectral interpretation of Riemann zeta zeros,
\newblock arXiv:2606.17494 (2026).

\bibitem{Hellinger1909}
E.~Hellinger,
\newblock Neue Begr\"undung der Theorie quadratischer Formen von unendlichvielen Ver\"anderlichen,
\newblock \emph{J. Reine Angew. Math.} \textbf{136} (1909), 210--271.

\bibitem{Kakutani1948}
S.~Kakutani,
\newblock On equivalence of infinite product measures,
\newblock \emph{Ann. of Math. (2)} \textbf{49} (1948), 214--224.

\bibitem{Meyer2005}
R.~Meyer,
\newblock A spectral interpretation for the zeros of the Riemann zeta function,
\newblock arXiv:math/0412277 (2005).

\bibitem{Nagoshi2016}
H.~Nagoshi,
\newblock The hypertranscendence of $L$-functions for
$\mathrm{GL}_m(\mathbb A_{\mathbb Q})$,
\newblock \emph{Bull. Aust. Math. Soc.} \textbf{93} (2016), 388--399.

\bibitem{Ostrowski1920}
A.~Ostrowski,
\newblock \"Uber Dirichletsche Reihen und algebraische Differentialgleichungen,
\newblock \emph{Math. Z.} \textbf{8} (1920), 241--298.

\bibitem{Pong2016}
W.-Y.~Pong,
\newblock Applications of differential algebra to algebraic independence of arithmetic
functions,
\newblock \emph{Acta Arith.} \textbf{172} (2016), 149--173.

\bibitem{Schwartz1966}
L.~Schwartz,
\newblock \emph{Th\'eorie des distributions},
\newblock Hermann, Paris, 1966.

\bibitem{Titchmarsh1986}
E.~C.~Titchmarsh,
\newblock \emph{The Theory of the Riemann Zeta-Function}, 2nd ed.,
\newblock revised by D.~R.~Heath-Brown, Oxford University Press, 1986.

\end{thebibliography}

\end{document}
